
\chapter{Literature Review}
\section{ Introduction}


Credit card companies must distinguish fraudulent from non-fraudulent transactions so that 
their customers' accounts will not get affected and charged for products they did not buy (Maniraj 
et al.,2019). Many financial Companies and institutions lose massive amounts of money because 
of Fraud and fraudsters that are seeking different approaches continuously to violate the rules and 
commit illegal actions; therefore, systems of fraud detection are essential for all banks that issue 
credit cards to decrease their losses (Zareapoor et al., 2012).
Multiple methods are used to detect 
fraudulent behaviours, such  as  Neural  Networks  (N.N.),  Decision  Trees,  K-nearest  neighbour 
algorithms, and Support Vector Machines (SVM). Those ML methods can be applied independently 
or collectively with ensemble or meta-learning techniques to develop classifiers (Zareapoor et al., 
2012).





\section{ Literature Review }
Zareapoor and his research team used multiple techniques to determine the best-performing model for detecting fraudulent transactions, which was established using the Accuracy of the model, the speed of detection and the cost. The models used were Neural Network, Bayesian Network, SVM, KNN and. The comparison table in theresearch paper showed that the Bayesian Network  was  high-speed in finding  fraudulent  transactions with  high  Accuracy.  The  N.N. performed well, as the detection was fast, with a medium accuracy. KNN’s speed was good with medium Accuracy, andfinally, SVM scored one of the lower scores, as the speed was low, and the Accuracy wasmedium. As for the cost, All models built were expensive (Zareapoor et al., 2012).\cite{1}\\


The model used by Alenzi and Aljehane to detect Fraud in credit cards was Logistic Regression. Their model scored 97.2\% in Accuracy, 97\% sensitivity and 2.8\% Error Rate. A comparison was performed between their model and two other classifiers ,  which are3 They were voting Classifier and KNN. V.C. scored 90\% in Accuracy, 88\% sensitivity and 10\% error rate, as for KNN where k = 1:10, the Accuracy of the model was 93\%, the sensitivity 94\%and 
7\% for the error rate (Alenzi & Aljehane, 2020).\cite{2}\\


Maniraj's team built a model to recognize if any new transaction is Fraud or non-fraud. Their 
goal was to get 100\% in detecting fraudulent transactions andtry to minimize the incorrectly classified fraud instances. Their model has performed well as they got 99.7\% of the fraudulent 
transactions (Maniraj et al., 2019).\cite{3}\\


The  classification  approach  used  by  Dheepa  and  Dhanapal  was  the  behaviour-based 
classification approach, using a Support Vector Machine, where the behavioural patternsof the customers were analyzed to distinguish credit card fraud, such as the amount, date,time, place, and frequency of card usage. The Accuracy achieved by their approach wasmore than 80\% (Dheepa & 
Dhanapal, 2012).\cite{4}\\


Mailini and Pushpa proposed using KNN and Outlier detection in identifying credit card fraud. After performing their model oversampled data, the authors found that the most suitable method for detecting and determining target instance anomaly is KNN, which showed that it is most suited to detecting Fraud with memory limitation. As for Outlier detection, the computation 
and memory required for credit card fraud detectionis much less in addition to working faster and better in large online datasets. However, their work and results showed that KNN was more accurate and efficient (Malini & Pushpa,2017).\cite{5}\\


Maes and his team proposed using Bayesian and Neural Networks to detect credit card fraud. 
Their results showed that Bayesian performance is 8\% more effective in detecting Fraud than 
ANN, meaning  BBN sometimes detects 8\% more fraudulent transactions. In addition to the 
Learning times, ANN can go up to several hours,whereas BBN takes only 20 minutes (Maes et al., 2002).\cite{6}\\


Jain’s team used several ML techniques to distinguish credit card fraud; three of them are SVM, ANN and KNN. Then, to compare the outcome of each model, they calculated the true positive (T.P.), false Negative (F.N.), false positive (F.P.), and true negative (T.N.) generated. ANN scored 99.71\% accuracy, 99.68\% precision, and 0.12\% false alarm rate.SVM accuracy is
94.65\%, 85.45\% for the precision, and 5.2\% false alarm rate. Moreover, finally,the Accuracy of KNN is 97.15\%, the precision is 96.84\%, and the false alarm rate is 2.88\% (Jain et al., 2019).\cite{7}\\


Dighe and his team used KNN, Logistic Regression and Neural Networks, multilayer perceptron and Decision Tree in their work, then evaluated the results regarding numerous accuracy metrics. Of all the models created, the best performing one is KNN, which scored 99.13\%, then in second place performing model at 96.40\% and in last place is logistic Regression 
with 96.27\% (Dighe et al., 2018).\cite{8}\\


Sahin and Duman used four Support Vector Machine methods in detecting credit card fraud. (SVM) Support Vector Machine with RBF, Polynomial, Sigmoid, and Linear Kernel,all models scored 99.87\% in the training model and 83.02\% in the testing part of the model (Sahin & Duman, 2011).\cite{9}\\

